% 05_results_B.tex — Results B: attribution / XAI. Drafted by item P2-E8-6 (P2 Sprint 7).
% Voice: ../STYLE.md (Andreas Maier). Numbers: tools/xai_study/compare/out/leaderboard.md
% + faithful_demo.md + the §R per-method records under tools/xai_study/phaseB_attribution/out/.
% Honesty contract: ../plan.md — faithfulness only; no semantics/meaning recovered here.

\subsection{Attribution: which inputs caused the output}\label{sec:results-b}

Attribution returns a map of which inputs mattered. We apply these methods to the VCS
to ask a question a real network cannot answer: is the map true?
Every attribution is scored against the same intervention oracle (\Sectionname\ref{sec:oracle})
that defines ground truth. For each chosen output $y$, a pixel, a score,
or a game event, we know the exact set of causes it should have named. We ran twelve
attribution methods across the 44 scored games, with the six core games (Pong, Breakout,
Space~Invaders, Seaquest, Ms.~Pac-Man, Q*bert) as running examples, plus an
oracle-as-method positive control, and score each against the
intervention oracle in Table~\ref{tab:phaseB}. We score all methods from a shared live
state rather than a boot frame. We seed a single random-action prefix of $90$ steps and a
horizon of $15$ (seed~$0$), and check the state with the oracle's cause-density gate: at
least $4$ candidate causes must move the position output (below) by more than $0.5$. The gate flags a state;
it does not reject it. All 44 scored games pass. Every method explains the same two outputs
per game (\Sectionname\ref{sec:methods}). The content output is one RAM byte at the end of
the horizon. The position
output is the number of screen pixels, over the whole frame, that differ from the
un-intervened frame; the gradient methods take its derivative through the bilinear sampler
at one sprite cell. The live state matters: on Seaquest the old boot-frame content oracle
moved only $2$ of $48$ candidate causes, while at the gameplay state $19$ of $48$ move the
output above the floor, so the methods are scored where the program is actually computing.

The main result is a split between two output types, not a ranking of methods. A
content output is the value of a register, a colour, or a graphics bit. It is a
differentiable function of the state, so its gradient flows. A position output is where
a sprite is drawn, set by strobe timing on a discrete index. It has provably zero naive
gradient (Paper~1, and \Sectionname\ref{sec:methods}). Hence the whole gradient family hits this
limit, and it splits the leaderboard into two clear groups.

% The former shared attribution/mechanistic figure (fig:fig4) is removed: per-phase results
% are now reported as tables --- attribution in Table~\ref{tab:phaseB}, mechanistic
% interpretability in Table~\ref{tab:phaseC} --- and the cross-tradition comparison is the
% leaderboard plane (Fig.~\ref{fig:faithfulness}) in \Sectionname\ref{sec:results_compare}.

On the content path the gradient is faithful, as expected. Vanilla
saliency~\cite{simonyan2014deep} puts most of its attribution on the true causal byte: on
Q*bert its attribution of a content read correlates with the oracle at Pearson
$0.999$ with precision@3 of $1.0$, on Ms.~Pac-Man at $0.811$ / $1.0$, on Breakout
at $0.750$ / $1.0$. Integrated Gradients (IG)~\cite{sundararajan2017axiomatic},
Grad$\times$Input/DeepLIFT~\cite{shrikumar2017learning},
SmoothGrad~\cite{smilkov2017smoothgrad}, and guided
backpropagation~\cite{springenberg2015striving} all reproduce this, because the
output we explain is a one-hot read whose Jacobian $\partial y/\partial u$ is a
constant pointing at the read index. Smoothing does not change the score
($\Delta\,\mathrm{corr}\approx 0$), and guided backprop reduces to vanilla
saliency. On this path all methods give the same answer.

The gradient family is weak on the position output. The naive gradient there is
identically zero, so the raw saliency, Grad$\times$Input, SmoothGrad, guided-backprop,
and IG maps carry no signal through the position read
(Prop.~\ref{prop:zero}). A differentiable substrate can restore a non-zero gradient here.
The battery scores the position output with that repair, the bilinear sampler, switched on
(\Sectionname\ref{sec:results_compare}), but the restored gradient
is faithful only on a minority of games and its faithfulness stays low. Averaged over the
position regime the gradient and correlational family reaches only $0.22$ against the
oracle, against $0.56$ for the causal and intervention family (Table~\ref{tab:phaseB}).
Across all output regimes the gradient and correlational family scores $0.38$ and the
causal and intervention family $0.71$. The position map is not empty by choice; it is
weak exactly where the discrete game logic sits.

One objection remains: these gradients differentiate a single read at the intervention
frame, not the program that runs during the horizon. We therefore ran a true end-to-end
control on the six core games. We read the output at the end of the horizon and
differentiate it back to the RAM at the intervention frame through every CPU instruction
in between, $65{,}081$ to $99{,}703$ instructions per game. The pass is checkpointed every
$200$ instructions, and the forward pass is bit-exact to the hard machine in every game.
TIA and RIOT reads enter as constants, so the renderer itself is not differentiated. The
result is the same one-hot map. In five of the six games the end-to-end gradient of the
content byte is exactly $e_{\mathrm{idx}}$ ($\lVert g\rVert_1=1$, one non-zero entry), and
its scores equal those of the single read to the last digit (Breakout $0.750$,
Space~Invaders $1.000$, Seaquest $0.017$, Ms.~Pac-Man $0.811$, Q*bert $0.999$). Every other
path by which the oracle's causes reach the output runs through a branch, an indexed load,
or a hardware read, and the straight-through gradient carries none of these. Pong is the
one exception, and it is instructive: the gradient also reaches \texttt{ram[56]} with
$\partial y/\partial u = 15$, the horizon length, which is what a per-frame velocity added
to a position produces. That byte is not among the candidate causes, so Pong's scores are
unchanged ($0.464$ content, $0.278$ position). On the position output the end-to-end
gradient vanishes entirely on Breakout and Space~Invaders, although the sampler's gradient
at the end of the horizon is non-zero. Breakout recomputes its paddle position every frame
from the paddle input, a hardware read the gradient cannot see. Differentiating the whole
program therefore does not rescue the gradient family; it sees the same smooth sliver of
the computation.

Two further results make this point stronger. The baseline choice is the perennial
worry about IG. Here it changes magnitude, not faithfulness. The
$(x-x_0)$ factor moves $\max|\mathrm{IG}|$, yet the ranking and the oracle correlation
hold, because the constant one-hot Jacobian keeps the mass on the same byte. Only the
degenerate baseline equal to the content byte breaks it.
Expected Gradients~\cite{erion2021improving} averages IG over a distribution of
recorded VCS states. It removes that degeneracy and is stable across reference draws: over five
reference seeds its oracle correlation varies by less than $10^{-15}$ on every scored game.
On $20$ of the $44$ games, however, its content attribution is identically zero, mostly
because the content byte holds the same value in every reference state, so every $(x-x_0)$
factor vanishes. On the content path it matches a single-baseline IG on $25$ games, falls below
it on $17$ ($16$ by more than $0.05$) and exceeds it on $2$ (by at most $0.04$), at
$K\times$ the compute (for $K$ baseline draws), and it still vanishes on position. The more careful method gives stability but not accuracy. Its
all-regime faithfulness of $0.19$ is the lowest among the attribution methods.

Guided backprop and vanilla saliency fail the program-randomization sanity check of
Adebayo et al.~\cite{adebayo2018sanity} on the system itself. We scramble the ROM so
the program boots to a completely different state (on average $70$--$127$ of $128$ RAM
bytes change, lowest on Seaquest), and the content-path saliency map stays byte-identical
(cosine $\approx 1.0$ on all six games). The Jacobian $\partial y/\partial u = e_{\mathrm{idx}}$ is a
constant one-hot, independent of what the program computes. So the map points at the
read index no matter which program produced the output. The saliency is therefore
provably model-invariant. A faithful explanation must change when the model changes.
This one does not. The intervention oracle, by contrast, distinguishes the two programs at
once. The map identifies the right cause but does not depend on the model that produced it.

When we replace the gradient with a real intervention, the position result improves.
Occlusion~\cite{zeiler2014visualizing} sets a candidate cause and re-runs
the bit-exact program, a coarse version of the oracle. On the position regime, where the
naive gradient is zero, it reaches faithfulness $0.619$, and its all-regime faithfulness
of $0.685$ is the highest among the attribution methods. The black-box perturb-and-re-run
methods follow the same logic and result:
KernelSHAP / Shapley sampling~\cite{lundberg2017unified,strumbelj2014explaining}
($0.650$), LIME~\cite{ribeiro2016why} ($0.640$), and RISE~\cite{petsiuk2018rise}
($0.544$) each draw masked coalitions (random held-out subsets of candidate causes) and average over the genuine responses, so each
works on position outputs where the gradient is dead, and each reports its own
completeness ($\Sigma\phi = y(\text{intact}) - y(\text{all-masked})$ for KernelSHAP,
LIME's surrogate $R^2$, the fit of its local linear model). The leaderboard nevertheless
tags these three as attribution-tradition methods and counts them in the gradient and
correlational family, not in the causal and intervention family.

The two minimal-mask methods meet Atrey et al.'s~\cite{atrey2020exploratory} objection
that perturbations stray off-manifold (a byte value the running program would never itself write). Extremal
perturbation~\cite{fong2017interpretable,fong2019understanding} optimises an
area-bounded mask of real $do(\text{occlude})$ re-runs, and the on-distribution
counterfactual~\cite{olson2021counterfactual} substitutes candidate cells toward a
genuine alternative VCS state, the RAM of another frame of the same ROM. Both stay on
the data manifold by construction and both work on position outputs (position-regime
faithfulness $0.34$ and $0.28$, all-regime $0.46$ and $0.43$). The
counterfactual has no effect (a no-op) on some games, for example where no on-distribution
content byte varies between frames. Its map is then flat, and these cells enter its mean
as $F=0$; we do not drop them. The cross-tradition contrast is clearest across all three phases: the causal and
intervention family averages $0.71$ faithfulness [$n=11$ methods]
against the oracle, against the gradient and correlational family's $0.38$
[$n=14$] (Table~\ref{tab:phaseB}). Restricted to the attribution methods
alone the two means are closer, because the content-path gradient is faithful and lifts the
gradient mean. The intervention advantage is a cross-tradition and position-regime
effect, not an attribution-wide one, and we report it as such. A perturbation is faithful
only when it is a valid intervention.

Some popular methods return no answer at all, because the VCS has no structure for
them to use. Grad-CAM and Grad-CAM++~\cite{selvaraju2017gradcam,%
chattopadhyay2018gradcampp} need a final convolutional layer to weight; attention
rollout~\cite{abnar2020quantifying} needs attention matrices;
VIPER~\cite{bastani2018verifiable} needs a learned policy to distil. The VCS is a
6507 CPU (a reduced-pin 6502) with a TIA and a RIOT executing a fixed ROM. Across all six games the
count of each required substrate is exactly zero, against $\geq 16$ genuine candidate
causes per game for the applicable methods. We record this as a measured finding,
$\mathtt{applies = false}$, not a poor score: the absence is structural and the same
in every game, and re-targeting Grad-CAM's pooling onto raw register state would only
be the content-path gradient renamed. Several well-known methods produce no result on
a real artifact. This absence is itself a finding of the audit.

Across the twelve methods, faithfulness and plausibility move in opposite directions:
methods rank high on plausibility but low on faithfulness. The clearest pair makes it
concrete. On the position regime the causal method activation patching (\Sectionname\ref{sec:results-c}) reaches
faithfulness $1.000$, the oracle ceiling, while vanilla saliency, a
widely used baseline, reaches only $0.092$. This is a per-method position gap of $0.908$.
Yet, on the plausibility proxy of \Sectionname\ref{sec:methods}
the popular method has the higher score ($0.9$ versus $0.5$): the weaker map has the
higher plausibility score. Low faithfulness with high plausibility
is the danger zone (\Sectionname\ref{sec:methods}). On the VCS we can measure it directly, not just warn about it.

We claim only that these maps are faithful, nothing more. A correct occlusion map names
the true causal bytes. It does not tell us what those bytes mean, that a cell is the
ball's $x$ or that a routine is the bounce. Object-level semantics enter only through
the external Tier-3 (T3) labels, game-concept names imported from AtariARI and OCAtari (\Sectionname\ref{sec:methods}), and even there we
verify a supplied label by intervention rather than discover it. Whether a faithful
attribution amounts to understanding the computation we take up in
\Sectionname\ref{sec:results_compare}. Here we have established only that faithfulness is
achievable, and that the most popular methods do not achieve it.

\begin{table}[t]
\footnotesize
\setlength{\tabcolsep}{4pt}
\caption{\textbf{Attribution scored against the intervention oracle
(\Sectionname\ref{sec:oracle}).} Twelve attribution methods aggregated over the 44 scored games
(oracle-as-method positive control $=1.0$). $F$ is the all-regime faithfulness against
ground truth with its $95\%$ bootstrap CI over games; $F_{\text{pos}}$ is the position/index regime
alone, where the naive gradient is provably zero; \emph{Plaus} is the documented
method-tradition plausibility \emph{proxy} (\Sectionname\ref{sec:methods}), \emph{not} a human
measurement. The naive gradient is provably zero on the position read
(Prop.~\ref{prop:zero}); the small non-zero $F_{\text{pos}}$ the gradient methods still
carry is the sampler-restored gradient of \Sectionname\ref{sec:results_compare}, with which
the battery scores every position output; it is faithful on only a minority of games and
stays low. Across all regimes the causal and intervention family
leads ($0.71$ against $0.38$, position $0.56$ against $0.22$), and the popular gradient
methods carry the \emph{higher} plausibility tag --- the danger zone, measured. The two
family-mean rows are the cross-tradition families ($n=14$ gradient/correlational, $n=11$
causal/intervention across all three traditions), not the twelve rows above; restricted
to attribution alone the two means are closer,
because the faithful content gradient lifts the gradient mean. In the \emph{Family}
column, \emph{interv.}\ rows belong to the causal/intervention family; \emph{gradient} and
\emph{black-box} (perturb-and-re-run) rows are counted in the gradient/correlational family,
following the leaderboard's tradition tags. The on-distribution
counterfactual is a no-op on some core games such as Space~Invaders, Ms.~Pac-Man, and
Q*bert (the substitution toward the alternative frame moves nothing, so $\Delta=0$, the map
is flat, and these cells enter its mean as $F=0$). Grad-CAM, Grad-CAM++, attention rollout, and
VIPER~\cite{selvaraju2017gradcam,abnar2020quantifying,bastani2018verifiable} do not apply
(no convolution, attention, or learned policy) and are omitted. Values read from the
committed leaderboard~\cite{leaderboardP2}.}\label{tab:phaseB}
\begin{tabularx}{\textwidth}{@{}X l l c c@{}}
\toprule
Method & Family & $F$ ($95\%$ CI) & $F_{\text{pos}}$ & Plaus \\
\midrule
Occlusion~\cite{zeiler2014visualizing}            & interv.\  & $0.69$ $[0.63,0.74]$ & $0.62$ & $0.55$ \\
KernelSHAP~\cite{lundberg2017unified}             & black-box & $0.65$ $[0.59,0.71]$ & $0.59$ & $0.90$ \\
LIME~\cite{ribeiro2016why}                        & black-box & $0.64$ $[0.58,0.70]$ & $0.58$ & $0.90$ \\
RISE~\cite{petsiuk2018rise}                        & black-box & $0.54$ $[0.50,0.59]$ & $0.43$ & $0.90$ \\
Grad$\times$Input / DeepLIFT~\cite{shrikumar2017learning} & gradient & $0.47$ $[0.44,0.51]$ & $0.09$ & $0.90$ \\
Extremal perturbation~\cite{fong2019understanding} & interv.\  & $0.46$ $[0.41,0.51]$ & $0.34$ & $0.55$ \\
On-distribution counterfactual~\cite{olson2021counterfactual} & interv.\  & $0.43$ $[0.35,0.51]$ & $0.28$ & $0.55$ \\
SmoothGrad~\cite{smilkov2017smoothgrad}           & gradient & $0.36$ $[0.30,0.42]$ & $0.09$ & $0.90$ \\
Vanilla saliency~\cite{simonyan2014deep}          & gradient & $0.36$ $[0.30,0.42]$ & $0.09$ & $0.90$ \\
Integrated Gradients~\cite{sundararajan2017axiomatic} & gradient & $0.36$ $[0.30,0.42]$ & $0.09$ & $0.90$ \\
Guided backprop~\cite{springenberg2015striving}   & gradient & $0.32$ $[0.26,0.37]$ & $0.00$ & $0.90$ \\
Expected Gradients~\cite{erion2021improving}      & gradient & $0.19$ $[0.13,0.25]$ & $0.03$ & $0.90$ \\
\midrule
\multicolumn{2}{@{}l}{Gradient/corr.\ mean}  & $0.38$ & $0.22$ & --- \\
\multicolumn{2}{@{}l}{Causal/interv.\ mean}  & $0.71$ & $0.56$ & --- \\
\botrule
\end{tabularx}
\end{table}

% =====================================================================
% BIBTEX-NEEDED: full entries for the keys this section cites that are
% NOT yet in references.bib (SM to merge; do NOT edit references.bib here).
% Already-present keys reused as-is (NO entry below): sundararajan2017axiomatic,
% selvaraju2017gradcam, chattopadhyay2018gradcampp, atrey2020exploratory.
% =====================================================================
%
% @inproceedings{simonyan2014deep,
%   title={Deep Inside Convolutional Networks: Visualising Image Classification Models and Saliency Maps},
%   author={Simonyan, Karen and Vedaldi, Andrea and Zisserman, Andrew},
%   booktitle={International Conference on Learning Representations (ICLR) Workshop},
%   year={2014}
% }
%
% @inproceedings{shrikumar2017learning,
%   title={Learning Important Features Through Propagating Activation Differences},
%   author={Shrikumar, Avanti and Greenside, Peyton and Kundaje, Anshul},
%   booktitle={Proceedings of the 34th International Conference on Machine Learning (ICML)},
%   pages={3145--3153},
%   year={2017}
% }
%
% @inproceedings{springenberg2015striving,
%   title={Striving for Simplicity: The All Convolutional Net},
%   author={Springenberg, Jost Tobias and Dosovitskiy, Alexey and Brox, Thomas and Riedmiller, Martin},
%   booktitle={International Conference on Learning Representations (ICLR) Workshop},
%   year={2015}
% }
%
% @inproceedings{smilkov2017smoothgrad,
%   title={SmoothGrad: Removing Noise by Adding Noise},
%   author={Smilkov, Daniel and Thorat, Nikhil and Kim, Been and Vi{\'e}gas, Fernanda and Wattenberg, Martin},
%   booktitle={ICML Workshop on Visualization for Deep Learning},
%   year={2017}
% }
%
% @article{erion2021improving,
%   title={Improving Performance of Deep Learning Models with Axiomatic Attribution Priors and Expected Gradients},
%   author={Erion, Gabriel and Janizek, Joseph D. and Sturmfels, Pascal and Lundberg, Scott M. and Lee, Su-In},
%   journal={Nature Machine Intelligence},
%   volume={3},
%   number={7},
%   pages={620--631},
%   year={2021},
%   publisher={Nature Publishing Group}
% }
%
% @inproceedings{zeiler2014visualizing,
%   title={Visualizing and Understanding Convolutional Networks},
%   author={Zeiler, Matthew D. and Fergus, Rob},
%   booktitle={European Conference on Computer Vision (ECCV)},
%   pages={818--833},
%   year={2014},
%   organization={Springer}
% }
%
% @inproceedings{fong2017interpretable,
%   title={Interpretable Explanations of Black Boxes by Meaningful Perturbation},
%   author={Fong, Ruth C. and Vedaldi, Andrea},
%   booktitle={Proceedings of the IEEE International Conference on Computer Vision (ICCV)},
%   pages={3429--3437},
%   year={2017}
% }
%
% @inproceedings{fong2019understanding,
%   title={Understanding Deep Networks via Extremal Perturbations and Smooth Masks},
%   author={Fong, Ruth and Patrick, Mandela and Vedaldi, Andrea},
%   booktitle={Proceedings of the IEEE/CVF International Conference on Computer Vision (ICCV)},
%   pages={2950--2958},
%   year={2019}
% }
%
% @inproceedings{petsiuk2018rise,
%   title={RISE: Randomized Input Sampling for Explanation of Black-box Models},
%   author={Petsiuk, Vitali and Das, Abir and Saenko, Kate},
%   booktitle={British Machine Vision Conference (BMVC)},
%   year={2018}
% }
%
% @inproceedings{ribeiro2016why,
%   title={``Why Should I Trust You?'': Explaining the Predictions of Any Classifier},
%   author={Ribeiro, Marco Tulio and Singh, Sameer and Guestrin, Carlos},
%   booktitle={Proceedings of the 22nd ACM SIGKDD International Conference on Knowledge Discovery and Data Mining},
%   pages={1135--1144},
%   year={2016}
% }
%
% @inproceedings{lundberg2017unified,
%   title={A Unified Approach to Interpreting Model Predictions},
%   author={Lundberg, Scott M. and Lee, Su-In},
%   booktitle={Advances in Neural Information Processing Systems (NeurIPS)},
%   volume={30},
%   pages={4765--4774},
%   year={2017}
% }
%
% @article{strumbelj2014explaining,
%   title={Explaining Prediction Models and Individual Predictions with Feature Contributions},
%   author={{\v{S}}trumbelj, Erik and Kononenko, Igor},
%   journal={Knowledge and Information Systems},
%   volume={41},
%   number={3},
%   pages={647--665},
%   year={2014},
%   publisher={Springer}
% }
%
% @article{olson2021counterfactual,
%   title={Counterfactual State Explanations for Reinforcement Learning Agents via Generative Deep Learning},
%   author={Olson, Matthew L. and Khanna, Roli and Neal, Lawrence and Li, Fuxin and Wong, Weng-Keen},
%   journal={Artificial Intelligence},
%   volume={295},
%   pages={103455},
%   year={2021},
%   publisher={Elsevier}
% }
%
% @inproceedings{adebayo2018sanity,
%   title={Sanity Checks for Saliency Maps},
%   author={Adebayo, Julius and Gilmer, Justin and Muelly, Michael and Goodfellow, Ian and Hardt, Moritz and Kim, Been},
%   booktitle={Advances in Neural Information Processing Systems (NeurIPS)},
%   volume={31},
%   pages={9505--9515},
%   year={2018}
% }
%
% @inproceedings{abnar2020quantifying,
%   title={Quantifying Attention Flow in Transformers},
%   author={Abnar, Samira and Zuidema, Willem},
%   booktitle={Proceedings of the 58th Annual Meeting of the Association for Computational Linguistics (ACL)},
%   pages={4190--4197},
%   year={2020}
% }
%
% @inproceedings{bastani2018verifiable,
%   title={Verifiable Reinforcement Learning via Policy Extraction},
%   author={Bastani, Osbert and Pu, Yewen and Solar-Lezama, Armando},
%   booktitle={Advances in Neural Information Processing Systems (NeurIPS)},
%   volume={31},
%   pages={2494--2504},
%   year={2018}
% }
